\section{Systems-Level Analysis \& Roofline Profiling}
\label{sec:systems}

\subsection{Roofline Model Formulation}
To rigorously quantify the hardware acceleration of NeuralOps, we apply the Roofline performance model. Let $\pi_{\text{peak}}$ denote the peak computational throughput of the tensor cores (in FLOPS), and let $\beta_{\text{peak}}$ denote the peak memory bandwidth of the HBM (in bytes/sec). The maximum achievable performance $P$ for an operation with arithmetic intensity $\mathcal{I}$ is given by:
\begin{equation}
P = \min(\pi_{\text{peak}}, \mathcal{I} \cdot \beta_{\text{peak}})
\end{equation}

By fusing the dual SVD projections and RMSNorm entirely within SRAM, NeuralOps shifts the operational point of linear attention from the memory-bound region ($\mathcal{I} < \pi_{\text{peak}} / \beta_{\text{peak}}$) into the compute-bound region, maximizing tensor core utilization.

\subsection{Kernel Operation Breakdown}
The SRAM-fused kernel bypasses the standard multi-kernel launch overhead. Specifically, it eliminates three roundtrips to HBM per attention block, translating to a $60\%$ reduction in global memory traffic and directly yielding the observed $2.5\times$ end-to-end acceleration.
